\section{Data and Lexicon}
\label{sec:data}

\subsection{The ViHSD corpus}
\label{sec:vihsd}

We use UIT-ViHSD \cite{luu2021vihsd}, a corpus of Vietnamese social-media comments
annotated with three mutually exclusive labels: \clean{} (no offensive content),
\off{} (offensive language that does not target a protected group as hate speech)
and \hate{} (hateful or seriously derogatory content). The \off{}/\hate{}
distinction depends on target and intensity rather than on a particular word, so a
purely lexical rule engine is insufficient on its own.

After cleaning, the corpus contains 33{,}398 comments. We use the fixed split shipped
with the project, and every system---rule engine, backbone, fusion variant and
analysis script---reads the same files. Each run directory stores the SHA-256 digest
of the split it used and re-checks it before training, ruling out a silently
re-shuffled split. Training contains 24{,}046 comments (2{,}556 \hate{}, 1{,}605 \off{}, 19{,}885
\clean{}), validation 2{,}672 (270/212/2{,}190) and test 6{,}680 (688/444/5{,}548).

The test split is imbalanced: 83.05\% \clean{}, 10.30\% \hate{} and 6.65\% \off{}.
Accuracy is therefore a poor headline metric---a constant \clean{} predictor already
reaches 0.8305---so macro-F1 is primary, always reported with per-class F1. The
6{,}680 test rows contain 6{,}571 distinct texts, i.e.\ 109 duplicated rows. They
are kept because removing them would change the frozen test set; one appears twice
among the 15 regressions, which therefore cover 14 distinct texts.

\subsection{The hate-speech lexicon}
\label{sec:lexicon}

The rule engine is driven by a Vietnamese hate-speech lexicon of \textbf{668
entries}, mined and labelled from the \emph{training} portions of three corpora:
ViHSD; R-ViHSD, a cleaned release that adds TDTU YouTube and voz-forum comments to
ViHSD \cite{bontoriTdtu}; and the code-mixed ViCM corpus \cite{nguyen2026vicm}. The
three training splits yield 1{,}062 candidate terms, of which 257, 665 and 140 are first
attested in ViHSD, ViCM and R-ViHSD; labelling reduces them to the entries kept below.
No comment from any evaluation split contributes to the dictionary, and each entry
records which corpora supported it (414 of the 668 occur in all three).
Each entry stores the surface term and a canonical form, a lexical label (\hate{} or
\off{}) with a per-label severity, an entry confidence estimated from the training
data, a context-required flag, and a review status (\emph{verified} or
\emph{auto}); both statuses may match, but the flag is retained in the evidence
record.

The lexicon is deliberately not a list of slurs. Most entries---551 of 668
(82.5\%)---are multi-word constructions such as \emph{mả cha}, \emph{con mẹ} or
\emph{ngu hơn}, and 524 (78.4\%) are context-required, because Vietnamese insults
frequently work by comparison and implication rather than by a single taboo word.
Only 74 entries carry an \off{} lexical label against 594 for \hate{}, which is why the
rule engine is weak on \off{} and fusion keeps that channel mostly neural. Three
properties make the lexicon worth keeping: it is \emph{editable} (a new variant needs
no retraining), \emph{deterministic} (the same comment yields the same spans) and
\emph{inspectable} (each score decomposes into a term, a match type and a context
factor).

\paragraph{Ethics of examples.}
All quoted comments are public ViHSD test comments. Because they contain profanity
and slurs, we mask the offending characters with \texttt{*} in every table;
evaluation always runs on the unmasked text. No user identifiers are reproduced, and
names that occasionally appear inside comments are part of the public corpus text.
